%% The following is a directive for TeXShop to indicate the main file
%%!TEX root = diss.tex


%========================================
\chapter{Groups for walks with small steps}
%========================================

Suppose that the random walk model has a step polynomial $S$ in the variables $x_1,\ldots, x_d$. We are interested in the bi-rational transformations that leave the kernel $K = 1 - tS$ intact. For some $x = x_k$ we have bi-rational transformations sending $g: x \mapsto y$ that preserse $K$. That is $K(y) = K(x)$. Expanding the definition we have
\begin{align}
		1- tS(y) &= 1- tS(x)\\
		S(y) &= S(x)\\
		S_{-1}\ol{x} + S_1 x & = S_{-1} \ol{y} + S_1 y
\end{align}
Because of this equation it's equally valid to speak of the birational transformations preserving $S$. Collecting terms and simplifying yields
\begin{align}
	(S_{-1}-S_1 xy)(x-y) = 0
\end{align}
There are cases here. If both $S_{-1},S_1$ are zero the equation is trivially valid and $g$ can be anything. When exactly one is zero then the equation only has the trivial solution $y = x$. When neither is zero we have the two solutions $y= x$ and $y = \ol{x}S_{-1}/S_1$ in which case we say that the transformation is non-trivial. That is we define $g_k = \ol{x}_k S_k^{-1}/S_k^1$ whenever we have a non-trivial transformation. Note that
\begin{align}
	g_k^2(x_k) = g_k\left({S_k^{-1}}{x_kS_k^1}\right) = x_k
\end{align}
so these transformations are involutions. We can define a group $G_S$ of the non-trivial transformations.
\begin{definition} Let $G = G_S$ be the group generated by the non-trivial transformations $g_k$.
\end{definition}
In the case where the transformation on $x$ is trivial there is something we can say about the random walk model. We use the operator $E = 1 - \ext{\ol{x}}$ to gain information on the generating function $Q$. Since $Q$ only has non-negative powers of $x$ in its expansion $EQ = Q$. However in the master functional equation we have
\begin{align}
	ESQ = SQ - \ol{x}(S_{-1}Q_1 + S_0 Q_0 + S_1 Q_{-1})= SQ - \ol{x}(S_{-1}Q_1 + S_0 Q_0)
\end{align}
We have several cases
\begin{itemize}
\item[1)] If $S_{-1} = 0$ and $S_1 = 0$, $S = S_0$ and therefore $ESQ = S_0EQ = S_0Q$ so the functional equation doesn't "see" the variable $x$.
\item[2)] If $S_{-1} = 0$ and $S_1 \neq 0$ then
\begin{align}
	ESQ = SQ - \ol{x}(S_{-1}Q_0 + S_0 Q_{-1}) = SQ
\end{align}
since $Q_{-1} = 0$ so the extraction doesn't see the variable $x$. This is equivalent to saying that even though we are working in the positive orthant we can relax the constraint for the $x$ to the whole space since there are no steps going back in that direction. The variable $x$ is still present in the functional equation. If we have no steps that stay in the $x = 0$ hyperplane then $x$ essentially acts as a parameter in the problem and we can ignore it. Otherwise we'll have
\begin{align}
	LESQ = LSQ &= (1+b[x^0])SQ \\&= SQ + b(S_{-1}Q_1+S_0Q_0+S_1Q_{-1}) \\&= SQ + bS_0Q_0 \neq SQ.
\end{align}
\item[3)] If $S_{-1}\neq 0$ but $S_1 = 0$ then $ESQ$ is non-trivial but since we cannot take any steps away from the $x = 0$ hyperplane we can eliminate the $S_{-1}$ term in the analysis of the problem.
\end{itemize}
In all the cases the random walk model behaves as if $S_{-1} = 0$. In the non-interacting version of the problem we can fully ignore the $x$ coordinate in the analysis since it acts like a parameter. The problem is essentially one of lower dimension. In the interacting case we have a bit extra because of the marking operator. In either case we'll define the dimension of the group to be the number of non-trivial generators.
\begin{definition} Let the \emph{dimension} of the model or the group be the number of non-trivial generators of $G_S$ in the defining representation.
\end{definition}
Without loss of generality we will speak of models set in $d$ dimensions that have a full set of generators. Firstly, we can easily show that there are infinitely many different group sizes that are manifested among the different dimensions. Let $S$ be a $(d-1)$-dimensional model with group $G$ of largest finite order (number of elements) for that dimension. Set this model in $d$ dimensions and define a new model $S' = S(x_d+\ol{x}_d)$. This model has a non-trivial generator $g_d$ that $S$ lacks. Therefore the group $G'$ of this new model satisfies
\begin{align}
	G' \cong G \times \langle g_d \rangle
\end{align}
and $|G'| = 2|G|$ and continuing in this manner allows us to write an infinite sequence of distinct groups. That is, not only are there infinitely many distinct groups but also, there is a group in $d$ dimensions not encountered before in $(d-1)$ dimensions (neither in order nor structure). For example the orders of the finite groups set in 2 dimensions are 1,2,4,6 and 8. So in 3 dimensions we have groups of order 1,2,4,6,8,12 and 16 at least. In fact Steve's paper shows that there are even more - of orders 24 and 48 too.

%========================================
\section{Coxeter Groups}
%========================================

This discussion would be short without delving into the subject of Coxeter groups. The $d$-dimensional groups defined above are Coxeter groups - a generalization of the dihedral groups encountered in dimension 2 random walk models. Coxeter groups are defined as
\begin{align}
	G = \langle r_1,\ldots, r_d |(r_ir_j)^{m_{i,j}} = 1\rangle
\end{align}
where $m_{i,i} = 1$ and $m_{i,j} \ge 2$ for all $i\neq j$, and $m_{i,j}$ could also be infinite. We can call these $r$'s reflections. When $m_{i,j} = 2$ the reflections commute. Also $m_{i,j} = m_{j,i}$. The \emph{rank} of $G$ is the number of generators of in this specific representation of $G$. Often $G$ will have more than one representation with different numbers of generators. So in our discussion of groups for random walks with small steps the rank of $G$ is the dimension of the model. We can draw a graph for each Coxeter group
\begin{itemize}
\item label vertices with the generators of $G$
\item connect the vertices $g_i$ and $g_j$ if $m_{i,j} \ge 3$
\item label the edge with $m_{i,j}$ if $m_{i,j}\ge 4$.
\end{itemize}
If the graph has different components, the Coxeter group factors as a direct product of Coxeter groups represented by the individual components. We can see this in terms of the step polynomial $S$. Suppose that $S$ factors as
\begin{align}
	S = T(x_1,\ldots,x_k)R(x_{k+1},\ldots,x_d)
\end{align}
For a fixed $1 \leq i \leq k$ we have
\begin{align}
	S_i^j = [x_i^j]S = R[x_i^j]T = R T_i^j
\end{align}
and therefore
\begin{align}
	g_i(x_i) = \ol{x}_i\frac{T_i^{-1}R}{T_i^1R} = \ol{x}_i\frac{T_i^{-1}}{T_i^1}
\end{align}
Similarly if $k+1\leq j\leq d$ then
\begin{align}
	g_j(x_j) = \ol{x}_j\frac{R_i^{-1}}{R_i^1}
\end{align}
Clearly
\begin{align}
	g_ig_j = g_jg_i
\end{align}
since they don't interfere with each other. Since $i$ and $j$ were arbitrary, this means that the generators from different sets commute and the Coxeter graph of the group factors into two pieces. We can write
\begin{align}
	G_S = G_T \times G_R
\end{align}
where the subscripts highlight which parts of $S$ they groups are generated by. This factoring also allows us to write the generating function $Q$ as a Hadamard product of lower-dimensional generating functions. We can call each component of the graph a prime model. We can say that a simple small component is one whose group doesn't split non-trivially (its graph has one component).

It's worth mentioning that it doesn't seem that all Coxeter groups are manifested in random walk models. Looking at the classification of finite Coxeter groups, those of rank 2 have orders 4, 6, 8, 12, 10 and 2$p$ (manifested in a $p$-gon). Namely the groups of orders 10,12 and 2$p$ (for most $p$) are not manifested in random walk models with small steps on the quarter plane. It appears that the $p$-gon symmetries are not present for most $p$ because of the crystallographic nature of the problem. It's quite possible that the Crystallographic restriction theorem will apply to the random walk model groups in all dimensions. However, as can be seen, a lot of random walk model groups in 2 dimensions have generators that don't map the lattice back onto itself in a simple way because the transformations are bi-rational. We can define classes of groups.

\begin{itemize}
\item a model is "nice" iff $g_i(x_i) \in \ze[x_1,\ol{x}_1\ldots,x_d,\ol{x}_d]$
\item a model is "crystallographic" iff $g_i(x_1^{n_1}\ldots x_d^{n_d}) = x_1^{m_1}\ldots x_d^{m_d}$ for some integers $n_j$ and $m_j$.
\end{itemize}
The names are suggestive because the generators in crystallographic random walk models map monomials to monomials and we can see if $G$ is crystallographic. The next step up are "nice" models where monomials are mapped to Laurent polynomials. The rest of the classes have a generator mapping some monomials to rational functions. Perhaps those first two types of models can be represented via reflections and will thus have Weyl group.

\endinput
